\documentclass[11pt]{article}
\usepackage{amsmath,amssymb,amsthm}

\newtheorem{theorem}{Theorem}
\newtheorem{corollary}[theorem]{Corollary}
\newtheorem{lemma}[theorem]{Lemma}

\newcommand{\Up}{U_+}

\title{Same-depth inverse minima across a reset atom followed by $1113$}
\author{}
\date{27 August 2026}

\begin{document}
\maketitle

For a positive odd integer $x$, let
\[
 \Up(x):=\frac{3x+1}{2^{v_2(3x+1)}}.
 \tag{1}
\]
If $w=(a_1,\ldots,a_L)$ is an exact accelerated exponent word, put
\[
 B(w):=\sum_{i=1}^L a_i,
 \qquad
 C(w):=\sum_{i=1}^L3^{L-i}2^{a_1+\cdots+a_{i-1}}.
 \tag{2}
\]
Thus
\[
 2^{B(w)}\Up^L(x)=3^Lx+C(w).
 \tag{3}
\]
On an interior deficit-four record tail, write
\[
 A:=1,\qquad B:=12,\qquad C:=1113,\qquad D:=1122.
 \tag{4}
\]

For a positive odd target $y$, define $\mathcal R_L(y;M)$ to be the set
of reverse exponent strings $(d_1,\ldots,d_L)$ of total mass at most $M$
which give a positive odd $L$-step predecessor of $y$.  Explicitly, put
\[
 t_0:=y,
 \qquad
 t_i:=\frac{2^{d_i}t_{i-1}-1}{3},
 \tag{5}
\]
and require every $t_i$ to be a positive odd integer.  The forward word
from $t_L$ to $y$ is $(d_L,\ldots,d_1)$.

\begin{lemma}[inverse-mass recursion for the two reset--$C$ blocks]
Let $r\geq0$.
\begin{enumerate}
 \item If $y$ is the endpoint of the exact word
 \[
  W_B(r):=12\,1^r1113,
  \tag{6}
 \]
 then
 \[
  \boxed{
  \mathcal R_{r+6}(y;r+9)
  =\{(3,1^{r+3},2,1)\}.}
  \tag{7}
 \]
 \item If $y$ is the endpoint of the exact word
 \[
  W_D(r):=1122\,1^r1113,
  \tag{8}
 \]
 then
 \[
  \boxed{
  \mathcal R_{r+8}(y;r+12)
  =\{(3,1^{r+3},2,2,1,1)\}.}
  \tag{9}
 \]
\end{enumerate}
In each case the displayed reverse string is the reversal of the original
word.  Consequently the original word is the unique same-depth inverse
word of minimum total exponent.
\end{lemma}

\begin{proof}
Let $q$ be the start of the word.  Direct composition of the four atom
maps gives
\[
\begin{aligned}
 2^{r+9}y
 &=3^{r+6}q+13\,3^{r+4}-2^{r+7}
 &&\text{for }W_B(r),\\
 2^{r+12}y
 &=3^{r+8}q+137\,3^{r+4}-2^{r+10}
 &&\text{for }W_D(r).
\end{aligned}
\tag{10}
\]
Hence the residue class of $y$ modulo $3^L$, where $L$ is the relevant
word length, is independent of $q$.

The reverse sets obey the finite recursion
\[
\begin{aligned}
 \mathcal R_{\ell+1}(t;M)
 =\bigcup_{\substack{a\geq1,\ a\leq M-\ell\\
                     3\mid 2^at-1}}
 \left\{(a,\delta):
 \delta\in
 \mathcal R_{\ell}\!\left(\frac{2^at-1}{3};M-a\right)
 \right\}.
\end{aligned}
\tag{11}
\]
Only the residue of $t$ modulo $3^{\ell+1}$ is used in this recursion.
Substitution of the two residue classes from (10), followed by induction
on $r$, gives the following complete list after the first $r+4$ reverse
steps.  The last two columns record the residue of the current state and
the depth/mass still available.
\[
\begin{array}{c|c|c|c|c}
\text{family}&\text{live reverse prefix}&\text{mass}
 &t_{r+4}\bmod3&\text{remaining depth/mass}\\ \hline
W_B(r)&1,2,1,1^r,3&r+7&1&2/2\\
      &3,1^{r+3}&r+6&1&2/3\\ \hline
W_D(r)&1,2,1,1^r,3&r+7&1&4/5\\
      &3,1^{r+3}&r+6&1&4/6\\
      &3,1^{r+2},3&r+8&2&4/4.
\end{array}
\tag{12}
\]
The base case is obtained by four applications of (11).  In the induction
step, the only admissible transition which leaves enough mass for the
remaining depth inserts one exponent $1$ immediately before the final
entry displayed in each prefix.  This proves (12) for every $r$.

For completeness, the only lower-mass competitor through the terminal
atom $C$ begins with the reverse string
\[
 1,2,1,1.
 \tag{13}
\]
It is the half-seed branch.  Pulling (13) through the following $r$
reverse exponent-one steps reaches $(p-1)/2$.  Both reset atoms satisfy
\[
 p+1\equiv2\pmod3,
 \tag{14}
\]
so $3\mid(p-1)/2$ and this branch has no further inverse step.

The first row for $W_B(r)$ has current residue $1$, so its next exponent
is even and at least $2$; no mass then remains for the final inverse step.
The second row has exactly the mass needed for the forced suffix $(2,1)$.

For $W_D(r)$, the first row takes a forced exponent $2$ to another state
of residue $1$; a second forced exponent $2$ then leaves only one unit of
mass for two remaining steps.  The third row takes a forced exponent $1$
to a multiple of $3$ and stops.  The middle row has the unique forced
suffix $(2,2,1,1)$.  This proves (7) and (9).
\end{proof}

\begin{theorem}[reset--$C$ same-depth inverse minimum]
Let $r\geq0$, let $q$ be the start of one of the exact words
$W_B(r)$ or $W_D(r)$, and put
\[
 L:=|W_B(r)|=r+6
 \quad\text{or}\quad
 L:=|W_D(r)|=r+8.
 \tag{15}
\]
If $y=\Up^L(q)$, then
\[
 \boxed{
 q=\min\{z\in2\mathbb Z+1:z>0,\ \Up^L(z)=y\}.}
 \tag{16}
\]
The minimizer is unique.
\end{theorem}

\begin{proof}
Let $v=(c_1,\ldots,c_L)$ be any exact word from a positive odd $z$ to
$y$, and put $B'=B(v)$.  Lemma~1 gives
\[
 B'\geq B_0,
 \qquad
 B_0=
 \begin{cases}
  r+9,&W_B(r),\\
  r+12,&W_D(r).
 \end{cases}
 \tag{17}
\]
For $W_B(r)$, the start of the record atom $B$ is congruent to $11$
modulo $16$.  The value $q=11$ would make its endpoint $13$, whose next
exponent is $3$, whereas the next atom in $W_B(r)$ starts with exponent
$1$.  Hence $q\geq27$.  For $W_D(r)$, the exact $D$ congruence is
$q\equiv71\pmod {128}$, so $q\geq71$.  In particular,
\[
 q>2^{B_0+1-L}
 =\begin{cases}16,&W_B(r),\\32,&W_D(r),\end{cases}
 \tag{18}
\]
and the original affine equation implies $y>(3/2)^L$.

If $B'=B_0$, the same lemma says that $v$ is the original word, so
$z=q$.

Assume that $B'\geq B_0+1$.  For the partial sums in (2), positivity of
the exponents gives
\[
 c_1+\cdots+c_{i-1}\leq B'-(L-i+1).
\]
Consequently
\[
\begin{aligned}
 C(v)
 &\leq2^{B'-L}
       \sum_{i=1}^L3^{L-i}2^{i-1}\\
 &=2^{B'-L}(3^L-2^L)
 <2^{B'-L}3^L.
\end{aligned}
\tag{19}
\]
Equations (3) and (19), together with $y>(3/2)^L$, imply
\[
 z>
 \frac{2^{B'}y}{3^L}-2^{B'-L}
 \geq
 \frac{2^{B_0+1}y}{3^L}-2^{B_0+1-L}.
 \tag{20}
\]
If $C_0$ is the carry of the original word, then
$2^{B_0}y=3^Lq+C_0$, and hence
\[
 z>2q+\frac{2C_0}{3^L}-2^{B_0+1-L}.
 \tag{21}
\]
Since $C_0>0$, equation (18) makes the right side of (21) strictly larger
than $q$.  This proves (16) and uniqueness.
\end{proof}

\begin{corollary}[one-reset cross-block descent is impossible]
Suppose that an interior record tail contains a block $BA^rC$ or
$DA^rC$.  No positive odd inverse branch of the same raw depth reaches
the terminal endpoint from a state smaller than the block's actual start.
Therefore a descent argument which crosses a reset blocker before a
$C$ half-coalescence must either change the inverse depth or span at least
two nontrivial record atoms.
\end{corollary}

\begin{proof}
The first assertion is Theorem~2.  An inverse word spanning precisely the
displayed block has the same raw depth, so the second assertion follows.
\end{proof}

\end{document}
